\section{Record, inverse, and finite-record results}\label{sec:records}

This section gives the record-side results in compressed form.  Each main
result is read under the constrained D-MAP record assumption
\Cref{ass:model-inference-contract-main}, with any extra regularity or
finite-record protocol stated locally.  Between those statements, all
probabilities below are probabilities of the stationary visible binary record
generated by the sampled-counter labelled kernels \(T_0,T_1\).
The killed entries \(T_1(R,R)=T_1(D,R)=0\) are used as assumptions, not as
conclusions; \Cref{thm:two-state-renewal-characterization-main} states exactly
when same-bin recovery leakage preserves the renewal law on which the inverse
conclusions depend.
Throughout this section \((A_n)_{n\in\mathbb Z}\) denotes that two-sided
stationary extension.
One-sided strings such as \(A_0,\ldots,A_{N-1}\) are finite restrictions used
only for observed-record statistics; when a proof conditions on a past event
or shifts a renewal epoch, it is using the same stationary two-sided process.

\subsection{Model and renewal law}

\begin{proposition}[Constrained D-MAP labelled kernels]\label{prop:instrument-kernels-main}
Fix \(\Gamma,\kappa_{\mathrm r},\Delta\tau>0\).  With hidden states ordered as
\((R,D)\), let
\begin{equation}\label{eq:t0-t1-main}
T_0=
\begin{pmatrix}
p&0\\ b&s
\end{pmatrix},
\qquad
T_1=
\begin{pmatrix}
0&1-p\\0&a
\end{pmatrix},
\end{equation}
where
\begin{equation}\label{eq:p-c-main}
p=e^{-\Gamma\Delta\tau},
\qquad
s=e^{-\kappa_{\mathrm r}\Delta\tau},
\end{equation}
and
\begin{equation}\label{eq:a-b-main}
a=\int_0^{\Delta\tau}\kappa_{\mathrm r}e^{-\kappa_{\mathrm r}r}
\bigl(1-e^{-\Gamma(\Delta\tau-r)}\bigr)\,dr,
\qquad
b=\int_0^{\Delta\tau}\kappa_{\mathrm r}e^{-\kappa_{\mathrm r}r}
e^{-\Gamma(\Delta\tau-r)}\,dr .
\end{equation}
For \(\kappa_{\mathrm r}\ne\Gamma\),
\[
a=\frac{\kappa_{\mathrm r}(1-p)-\Gamma(1-s)}
        {\kappa_{\mathrm r}-\Gamma},
\qquad
b=\frac{\kappa_{\mathrm r}(p-s)}
        {\kappa_{\mathrm r}-\Gamma},
\]
with continuous equal-rate limits.  The labelled-kernel convention is
\((T_i)_{XY}=\PP(A_n=i,H_{n+1}=Y\mid H_n=X)\), and word probabilities are
\(\pi T_{u_0}\cdots T_{u_{m-1}}\mathbf1\) from the stationary hidden law.
These are the kernels fixed in \Cref{ass:model-inference-contract-main}.
\end{proposition}

\begin{proof}
The two rows define a labelled transition kernel directly.  From state \(R\),
the zero-labelled transition remains in \(R\) with probability \(p\), while the
one-labelled transition goes to \(D\) with probability \(1-p\).  From state
\(D\), the zero-labelled transition has weights \(b\) to \(R\) and \(s\) to
\(D\), and the one-labelled transition has weight \(a\) to \(D\).  The
integrals define nonnegative weights with \(a+b=1-s\), so the two row sums of
\(T_0+T_1\) are one.  Thus \(T_0,T_1\) form the stated D-MAP kernels.  The
displayed unequal-rate formulas are the elementary evaluations of the two
integrals, with continuous limits on the equal-rate diagonal.
\end{proof}

\begin{lemma}[Equal-rate and critical notation]\label{lem:equal-rate-critical-convention-main}
The survival factors are always
\[
p=e^{-\Gamma\Delta\tau},\qquad
s=e^{-\kappa_{\mathrm r}\Delta\tau},
\]
and throughout the stipulated-kernel algebra
\[
q=1-p .
\]
The hidden eigenvalue is
\begin{equation}\label{eq:lambda-hid-main}
\lambda_{\mathrm{hid}}
=p-b
=\frac{\kappa_{\mathrm r}s-\Gamma p}{\kappa_{\mathrm r}-\Gamma}
\quad(\kappa_{\mathrm r}\ne\Gamma),
\end{equation}
with continuous equal-rate value
\(\lambda_{\mathrm{hid}}=e^{-y}(1-y)\), where
\(y=\kappa_{\mathrm r}\Delta\tau=\Gamma\Delta\tau\) on the equal-rate
diagonal.  Put
\[
x=\Gamma\Delta\tau,\qquad y=\kappa_{\mathrm r}\Delta\tau .
\]
For \(x\ne y\),
\[
\lambda_{\mathrm{hid}}=0
\quad\Longleftrightarrow\quad
x e^{-x}=y e^{-y}.
\]
The equation \(x e^{-x}=y e^{-y}\) always has the formal solution \(x=y\).
For \(y\ne1\) this diagonal solution is only the removable zero of the
multiplied quotient numerator, not a zero of the continuous eigenvalue
\(\lambda_{\mathrm{hid}}\).  The critical rate is therefore the non-removable
branch
\begin{equation}\label{eq:gamma-crit-piecewise-main}
\Gamma_{\mathrm{crit}}(\kappa_{\mathrm r},\Delta\tau)
=\Delta\tau^{-1}
\begin{cases}
x_+(y),&0<y<1,\quad x_+(y)>1,\quad x_+(y)e^{-x_+(y)}=y e^{-y},\\
1,&y=1,\\
x_-(y),&y>1,\quad 0<x_-(y)<1,\quad x_-(y)e^{-x_-(y)}=y e^{-y}.
\end{cases}
\end{equation}
Equivalently, for \(y\ne1\) it is the unique positive solution of
\(x e^{-x}=y e^{-y}\) on the opposite side of \(1\) from \(y\), and the
removable equal-rate branch \(x=y\) is explicitly excluded.  Only at \(y=1\)
does the equal-rate branch become a genuine zero, giving
\(\Gamma_{\mathrm{crit}}=\kappa_{\mathrm r}=1/\Delta\tau\).
\end{lemma}

\begin{proof}
The identity \(p-b\) and the quotient form follow from
\eqref{eq:a-b-main}.  The equal-rate value is l'Hopital's rule, equivalently
\(b=y e^{-y}\) when \(\kappa_{\mathrm r}=\Gamma\), and hence
\(\lambda_{\mathrm{hid}}=e^{-y}(1-y)\).  For \(x\ne y\), multiplying the
quotient numerator by \(\Delta\tau\) gives the scalar equation
\(x e^{-x}=y e^{-y}\).  The function \(x\mapsto xe^{-x}\) is strictly
increasing on \((0,1)\), strictly decreasing on \((1,\infty)\), and has its
only maximum at \(1\).  Hence the level through \(y\ne1\) has exactly one
positive solution on the opposite side of \(1\), in addition to the removable
solution \(x=y\).  This gives the two non-removable branches in
\eqref{eq:gamma-crit-piecewise-main}.  When \(y=1\), the opposite branch
coalesces with the diagonal and the equal-rate limit \(e^{-1}(1-1)=0\) gives
\(\Gamma_{\mathrm{crit}}=1/\Delta\tau=\kappa_{\mathrm r}\).
\end{proof}

For later audits, the equal-rate diagonal
\(\Gamma=\kappa_{\mathrm r}\), \(y=\Gamma\Delta\tau\), \(p=e^{-y}\), is
summarized by the single dictionary
\begin{equation}\label{eq:equal-rate-dictionary-main}
\begin{aligned}
a&=1-(1+y)p,&
b&=yp,\\
\lambda_{\mathrm{hid}}&=p-b=(1-y)p,&
S_k&=\PP(G\ge k)=p^k(1+yk),\\
g_k&=p^k\bigl((1-p)(1+yk)-yp\bigr),&
h_k&=\frac{(1-p)+y(k(1-p)-p)}{1+yk},
\qquad k\ge0 .
\end{aligned}
\end{equation}

\begin{lemma}[Regeneration and zero-run law]\label{lem:renewal-gap-main}
Let \(G\) be the number of zeros between consecutive clicks in the stationary
constrained D-MAP record.  Click times are regeneration times.  For
\(k\ge0\),
\begin{equation}\label{eq:gap-main}
g_k=\PP(G=k)=e_D^\top T_0^kT_1\mathbf1 .
\end{equation}
If \(\kappa_{\mathrm r}\ne\Gamma\), then
\begin{equation}\label{eq:gap-closed-main}
g_k=
\frac{\kappa_{\mathrm r}(1-p)p^k-\Gamma(1-s)s^k}
     {\kappa_{\mathrm r}-\Gamma}.
\end{equation}
The gap law has full support on \(\mathbb N_0\) and an exponential tail.
\end{lemma}

\begin{proof}
After a click the next hidden state is \(D\), independently of the earlier
record.  Thus a gap of \(k\) zeros followed by a click has probability
\eqref{eq:gap-main}.  Since \(T_0\) is lower triangular,
\[
\textstyle\sum_{j=0}^{k-1}s^jp^{k-1-j}=0\quad(k=0),
\]
with the displayed sum interpreted as empty when \(k=0\), and
\[
T_0^k=
\begin{pmatrix}
p^k&0\\
b\textstyle\sum_{j=0}^{k-1}s^jp^{k-1-j}&s^k
\end{pmatrix}.
\]
Therefore
\[
e_D^\top T_0^kT_1\mathbf1
=b(1-p)\textstyle\sum_{j=0}^{k-1}s^jp^{k-1-j}+as^k,
\]
where the comma only separates the display from the following substitution.
For \(p\ne s\),
\[
\sum_{j=0}^{k-1}s^jp^{k-1-j}=\frac{p^k-s^k}{p-s}.
\]
Substituting
\[
a=\frac{\kappa_{\mathrm r}(1-p)-\Gamma(1-s)}
        {\kappa_{\mathrm r}-\Gamma},
\qquad
b=\frac{\kappa_{\mathrm r}(p-s)}
        {\kappa_{\mathrm r}-\Gamma}
\]
therefore gives
\[
b(1-p)\frac{p^k-s^k}{p-s}+as^k
=
\frac{\kappa_{\mathrm r}(1-p)p^k-\Gamma(1-s)s^k}
     {\kappa_{\mathrm r}-\Gamma},
\]
which is \eqref{eq:gap-closed-main}.  If \(p=s\), equivalently
\(\kappa_{\mathrm r}=\Gamma\), the same identity is read by the continuous
equal-rate dictionary \eqref{eq:equal-rate-dictionary-main}.  All terms are
probabilities of disjoint paths, so they are positive, and the two survival
factors lie in \((0,1)\).
\end{proof}

\begin{corollary}[Support of finite binary words]\label{cor:word-support-main}
Every finite binary word has positive probability under the stationary
constrained D-MAP record.
\end{corollary}

\begin{proof}
Let the word be \(w=(w_0,\ldots,w_{m-1})\), with \(m\ge1\), and write
\[
S=\{i\in\{0,\ldots,m-1\}:w_i=1\}.
\]
We work in the two-sided stationary renewal click set
\(\varXi=\{n:A_n=1\}\).  Embed the observed window in the longer interval
\([-1,m]\) and force clicks at both artificial endpoints.  More explicitly,
list the finite set
\[
C=\{-1\}\cup S\cup\{m\}=\{c_0<c_1<\cdots<c_J\}.
\]
Consider the event that \(c_0,c_1,\ldots,c_J\) are consecutive renewal
clicks, with no further clicks between them.  Under the stationary law this
event has probability
\[
\PP(c_0\in\varXi)\prod_{j=0}^{J-1}
\PP\bigl(G=c_{j+1}-c_j-1\bigr)
=\rho\prod_{j=0}^{J-1}g_{c_{j+1}-c_j-1}.
\]
Each displayed gap length is a nonnegative integer, including at the two
observed-window boundaries.  By \Cref{lem:renewal-gap-main},
\(\rho>0\) and \(g_k>0\) for every \(k\ge0\), so the event has positive
probability.  On this event the restriction of the click set to
\(\{0,\ldots,m-1\}\) is exactly \(S\): the forced interior clicks are precisely
the one-sites of \(w\), and the ``consecutive clicks'' condition rules out
clicks at every zero-site.  Hence the cylinder probability of \(w\) is at
least this positive probability.  The empty word has probability \(1\).
\end{proof}

\subsection{Exact dependence, hazard, and count diagnostics}

These are falsification diagnostics only; no size, power, or post-selection
coverage theorem is proved for these diagnostics.
The hazard, Fano, and low-lag formulas in this subsection have the inferential
status fixed in \Cref{ass:model-inference-contract-main}.  The hazard,
spectral-density, and Fano displays in this subsection are population
identities or inputs to later sandwich/projection procedures.

\begin{lemma}[Stationary hidden law]\label{lem:stationary-main}
For \(P=T_0+T_1=\begin{psmallmatrix}p&q\\ b&1-b\end{psmallmatrix}\), the
stationary hidden law is
\[
\pi=\left(\frac b{q+b},\frac{q}{q+b}\right).
\]
With \(q=1-p\) as fixed in \Cref{lem:equal-rate-critical-convention-main}, the
stationary click density is
\begin{equation}\label{eq:rho-main}
\rho=\pi T_1\mathbf1=\frac{q(1-s)}{q+b}.
\end{equation}
Moreover \(a<q(1-s)<q\), and therefore
\(a-\rho=-b(q-a)/(q+b)<0\).
\end{lemma}

\begin{proof}
Solving \(\pi P=\pi\) gives the displayed law.  Since
\(u=T_1\mathbf1=(q,a)^\top\) and \(a+b=1-s\), \(\rho=\pi u\) gives
\eqref{eq:rho-main}.  For \(r>0\),
\[
1-e^{-\Gamma(\Delta\tau-r)}<1-e^{-\Gamma\Delta\tau}=q,
\]
so the defining integral for \(a\) gives
\[
a<q\int_0^{\Delta\tau}\kappa_{\mathrm r}e^{-\kappa_{\mathrm r}r}\,dr
=q(1-s)<q .
\]
The sign follows from
\[
a-\rho=-\frac{b(q-a)}{q+b}
\]
and the strict positivity of \(b\) and \(q-a\).
\end{proof}

\begin{lemma}[Covariance transfer through the hidden mode]\label{lem:covariance-transfer-main}
For the stationary D-MAP record,
\begin{equation}\label{eq:cov-transfer-main}
\PP(A_n=1\mid A_0=1)
=\rho+(a-\rho)\lambda_{\mathrm{hid}}^{\,n-1},
\qquad
\operatorname{Cov}(A_0,A_n)=\rho(a-\rho)\lambda_{\mathrm{hid}}^{\,n-1}
\end{equation}
for \(n\ge1\).
\end{lemma}

\begin{proof}
With \(q=1-p\) and \(u=(q,a)^\top\),
\[
u-\rho\mathbf1=\frac{q-a}{q+b}\binom {q}{-b},
\qquad
P\binom {q}{-b}=(p-b)\binom {q}{-b}.
\]
Thus
\[
P^{n-1}u=\rho\mathbf1+
\lambda_{\mathrm{hid}}^{n-1}(u-\rho\mathbf1).
\]
A click at time \(0\) sets the next hidden state to \(D\).  Hence
\[
\PP(A_n=1\mid A_0=1)
=e_D^\top P^{n-1}u
=\rho+\lambda_{\mathrm{hid}}^{n-1}(a-\rho),
\]
because the \(D\)-component of \(u-\rho\mathbf1\) is \(a-\rho\).  Multiplying
by \(\rho=\PP(A_0=1)\) and subtracting \(\rho^2\) gives the covariance in
\eqref{eq:cov-transfer-main}.
\end{proof}

\begin{lemma}[Critical-rate uniqueness and covariance phase]\label{lem:critical-rate-main}
For fixed \(\kappa_{\mathrm r},\Delta\tau>0\), the equation
\(\lambda_{\mathrm{hid}}(\Gamma)=0\) has the unique positive solution
\(\Gamma_{\mathrm{crit}}(\kappa_{\mathrm r},\Delta\tau)\) defined by the
piecewise non-removable branch
\eqref{eq:gamma-crit-piecewise-main}.  When
\(\Gamma_{\mathrm{crit}}\ne\kappa_{\mathrm r}\), this is equivalently
\begin{equation}\label{eq:gamma-c-main}
\Gamma_{\mathrm{crit}}e^{-\Gamma_{\mathrm{crit}}\Delta\tau}
=\kappa_{\mathrm r}e^{-\kappa_{\mathrm r}\Delta\tau}.
\end{equation}
The displayed scalar equation also has the diagonal solution
\(\Gamma_{\mathrm{crit}}=\kappa_{\mathrm r}\), but that solution is removable
for \(\lambda_{\mathrm{hid}}\) unless
\(\kappa_{\mathrm r}\Delta\tau=1\).  Thus on the diagonal
\(\Gamma=\kappa_{\mathrm r}\), the zero occurs only at
\(\Gamma=\kappa_{\mathrm r}=1/\Delta\tau\).
The sign of \(\operatorname{Cov}(A_0,A_n)\) is the sign of
\(-\lambda_{\mathrm{hid}}^{\,n-1}\).
\end{lemma}

\begin{proof}
The uniqueness is the branch statement in
\Cref{lem:equal-rate-critical-convention-main}.  The covariance sign follows
from \(\rho>0\) and \(a-\rho<0\) in \Cref{lem:stationary-main}.
\end{proof}

\begin{lemma}[Hidden-mode contraction]\label{lem:hidden-mode-contraction-main}
For every positive constrained D-MAP setting,
\[
|\lambda_{\mathrm{hid}}|<1 .
\]
In particular the shifted geometric Fourier series in
\Cref{prop:spectrum-main} is absolutely convergent.
\end{lemma}

\begin{proof}
The transition matrix
\[
P=T_0+T_1=
\begin{pmatrix}
p&1-p\\ b&1-b
\end{pmatrix}
\]
is a two-state stochastic matrix with strictly positive off-diagonal entries
\(1-p>0\) and \(b>0\).  Its eigenvalues are \(1\) and
\(\lambda_{\mathrm{hid}}=p-b\).  The upper bound is immediate from \(p<1\)
and \(b>0\): \(\lambda_{\mathrm{hid}}<1\).  For the lower bound, \(p>0\) and
\[
b=\int_0^{\Delta\tau}\kappa_{\mathrm r}e^{-\kappa_{\mathrm r}r}
e^{-\Gamma(\Delta\tau-r)}\,dr
<\int_0^{\Delta\tau}\kappa_{\mathrm r}e^{-\kappa_{\mathrm r}r}\,dr
=1-s<1 .
\]
Thus \(\lambda_{\mathrm{hid}}=p-b>-1\).  Hence
\(|\lambda_{\mathrm{hid}}|<1\).  Absolute convergence follows from
\(\sum_{n\ge1}|\lambda_{\mathrm{hid}}|^{n-1}<\infty\).
\end{proof}

\begin{lemma}[Finite hidden-start path-space contraction]\label{lem:finite-hidden-start-path-contraction-main}
Fix one positive constrained D-MAP setting and put \(P=T_0+T_1\).  Let
\(\pi\) be the stationary hidden law.  For a hidden initial law \(\nu\), let
\(\PP_\nu\) denote the one-sided hidden/visible law started from \(H_0\sim\nu\),
and let \(\PP_\pi\) denote the stationary-start law.  There are constants
\(C<\infty\) and \(\theta\in(0,1)\), depending only on the fixed setting, such
that for every burn-in length \(L\ge0\), every retained length \(N\ge1\), and
every set \(B\subset\{0,1\}^N\),
\begin{equation}\label{eq:finite-hidden-start-path-tv-main}
\left|
\PP_\nu\bigl((A_L,\ldots,A_{L+N-1})\in B\bigr)
-
\PP_\pi\bigl((A_0,\ldots,A_{N-1})\in B\bigr)
\right|
\le C\theta^L .
\end{equation}
Equivalently,
\[
\sup_{B\subset\{0,1\}^N}
\left|
\PP_\nu\bigl((A_L,\ldots,A_{L+N-1})\in B\bigr)
-
\PP_\pi\bigl((A_0,\ldots,A_{N-1})\in B\bigr)
\right|
\le C\theta^L ,
\]
up to the harmless convention-dependent factor in the definition of total
variation.  The bound is uniform in \(N\).  Consequently, if \(L_N\to\infty\)
and \(S_N\) is any measurable statistic of the retained visible word
\((A_{L_N},\ldots,A_{L_N+N-1})\), then every stationary-start convergence in
distribution, convergence in probability, or event-probability limit for
\(S_N\) transfers unchanged to the finite hidden start.
\end{lemma}

\begin{proof}
For a word \(w=(w_0,\ldots,w_{N-1})\), write
\(T_w=T_{w_0}\cdots T_{w_{N-1}}\).  For a set
\(B\subset\{0,1\}^N\), put \(M_B=\sum_{w\in B}T_w\).  Since
\(\sum_{w\in\{0,1\}^N}T_w=P^N\), the matrix \(M_B\) is substochastic: its
entries are nonnegative and \(0\le M_B\mathbf1\le\mathbf1\) coordinatewise.
The retained visible word probability from the finite start is therefore
\[
\PP_\nu\bigl((A_L,\ldots,A_{L+N-1})\in B\bigr)
=\nu P^L M_B\mathbf1,
\]
while stationarity gives
\[
\PP_\pi\bigl((A_0,\ldots,A_{N-1})\in B\bigr)
=\pi M_B\mathbf1 .
\]
Thus the difference is
\[
(\nu P^L-\pi)M_B\mathbf1 .
\]
Because \(0\le M_B\mathbf1\le\mathbf1\), the total-variation dual inequality
gives
\[
\left|(\nu P^L-\pi)M_B\mathbf1\right|
\le \|\nu P^L-\pi\|_{\mathrm{TV}} .
\]

It remains only to bound the hidden-chain term.  The matrix
\[
P=\begin{pmatrix}p&1-p\\ b&1-b\end{pmatrix}
\]
has strictly positive off-diagonal entries, stationary law
\(\pi=(b/(1-p+b),(1-p)/(1-p+b))\), and nontrivial eigenvalue
\(\lambda_{\mathrm{hid}}=p-b\).  By
\Cref{lem:hidden-mode-contraction-main}, \(|\lambda_{\mathrm{hid}}|<1\).
On the two-state zero-sum subspace, multiplication by \(P\) is multiplication
by \(\lambda_{\mathrm{hid}}\).  Hence, for every initial law \(\nu\),
\[
\|\nu P^L-\pi\|_{\mathrm{TV}}
\le C|\lambda_{\mathrm{hid}}|^L
\]
with a finite constant \(C\) uniform over \(\nu\).  Taking any
\(\theta\in(|\lambda_{\mathrm{hid}}|,1)\), or \(\theta=|\lambda_{\mathrm{hid}}|\)
when \(\lambda_{\mathrm{hid}}\ne0\) and increasing \(C\) if desired, proves
\eqref{eq:finite-hidden-start-path-tv-main}.

The final assertion follows because every statistic of the retained word is a
measurable image of that \(N\)-word.  Total variation cannot increase under a
measurable map, so the law of \(S_N\) under the finite hidden start differs
from its stationary-start law by at most \(C\theta^{L_N}\to0\).  Applying this
bound to continuity sets gives distributional convergence; applying it to the
events \(\{|S_N-s|>\varepsilon\}\), or to the events appearing in a stated
probability limit, gives convergence in probability and event-probability
limits.
\end{proof}

\begin{lemma}[Hazard obstruction to finite symbolic memory]\label{lem:hazard-obstruction-main}
A renewal binary process with full gap support is a finite-order Markov chain
only if its gap hazard is eventually constant.
\end{lemma}

\begin{proof}
For a renewal process, after observing \(k\) consecutive zeros since the last
click, the conditional click probability is the zero-run hazard \(h_k\).
If an \(r\)-step Markov rule generated the same full-support law, then all
histories with \(k\ge r\) trailing zeros would have the same last \(r\) symbols
and hence the same click probability.  Therefore \(h_k\) would be constant for
all \(k\ge r\).

This condition concerns a transition rule conditioned on the most recent
symbols.  It is different from finite dependence of separated sigma-fields:
\(1\)-dependence says that events separated by one unobserved sample
factorize, while an \(r\)-step Markov rule says that every future one-step
transition probability is a function of only the last \(r\) symbols.  For a
renewal record with full gap support, the latter would force the click
probability after a long zero run to stop depending on the run length.
\end{proof}

\begin{lemma}[Cylinder factorization at the memory threshold]\label{lem:kernel-factorization-main}
If \(\lambda_{\mathrm{hid}}=0\), then \(P=T_0+T_1=\mathbf1\pi\).  More
explicitly, let \(u=(u_0,\ldots,u_{r-1})\) and
\(v=(v_0,\ldots,v_{\ell-1})\) be arbitrary finite binary words and write
\(T_u=T_{u_0}\cdots T_{u_{r-1}}\), \(T_v=T_{v_0}\cdots T_{v_{\ell-1}}\).
If the word \(u\) is observed on one block and \(v\) is observed on a second
block beginning two samples later, with the intervening sample unobserved,
then
\[
\PP(A_0^{r-1}=u,\ A_{r+1}^{r+\ell}=v)
=\pi T_u P T_v\mathbf1
=\bigl(\pi T_u\mathbf1\bigr)\bigl(\pi T_v\mathbf1\bigr).
\]
Consequently every two visible blocks separated by one unobserved sample are
independent, so the visible record is exact \(1\)-dependent.
\end{lemma}

\begin{proof}
The second eigenvalue of \(P\) is \(\lambda_{\mathrm{hid}}\), while the first is
\(1\).  At \(\lambda_{\mathrm{hid}}=0\), the primitive stochastic matrix
\(P\) has rank one and both rows equal \(\pi\), hence \(P=\mathbf1\pi\).
For fixed words \(u\) and \(v\), summing over the unobserved sample between
the two displayed blocks inserts \(T_0+T_1=P\).  The stationary cylinder
probability is therefore
\[
\pi T_u(T_0+T_1)T_v\mathbf1=\pi T_uPT_v\mathbf1.
\]
At the threshold,
\[
\pi T_uPT_v\mathbf1
=\pi T_u\mathbf1\,\pi T_v\mathbf1
=\PP(A_0^{r-1}=u)\PP(A_0^{\ell-1}=v),
\]
where the last equality is stationarity applied to the two single-block
cylinders.  Thus every pair of finite cylinder events supported on two
finite blocks with one omitted sample between them factorizes.  Translating
the indices gives the same statement for any blocks
\(\{i,\ldots,m\}\) and \(\{m+2,\ldots,j\}\).  Since finite cylinders generate
the corresponding finite-block sigma algebras, the two finite-block sigma
algebras are independent.  Letting the left and right block lengths increase
and using the monotone-class theorem gives independence of
\(\sigma(A_j:j\le m)\) and \(\sigma(A_j:j\ge m+2)\) for every \(m\).
\end{proof}

\begin{lemma}[Exact \texorpdfstring{\(1\)}{1}-dependence criterion]\label{lem:one-dep-criterion-main}
The constrained D-MAP record is exact \(1\)-dependent iff
\(\lambda_{\mathrm{hid}}=0\), equivalently
\(\Gamma=\Gamma_{\mathrm{crit}}(\kappa_{\mathrm r},\Delta\tau)\).
\end{lemma}

\begin{proof}
The forward direction at the threshold is
\Cref{lem:kernel-factorization-main}.  Conversely, exact \(1\)-dependence
would make \(A_0\) and \(A_2\) independent.  But
\[
\operatorname{Cov}(A_0,A_2)=\rho(a-\rho)\lambda_{\mathrm{hid}},
\]
and the prefactor is nonzero by \Cref{lem:stationary-main}.
\end{proof}

\begin{lemma}[Survival tails and hazard monotonicity]\label{lem:survival-hazard-main}
For \(\kappa_{\mathrm r}\ne\Gamma\),
\begin{equation}\label{eq:survival-tail-main}
S_k=\PP(G\ge k)=
\frac{\kappa_{\mathrm r}p^k-\Gamma s^k}{\kappa_{\mathrm r}-\Gamma},
\end{equation}
and
\begin{equation}\label{eq:hazard-main}
h_k=\PP(G=k\mid G\ge k)=
\frac{\kappa_{\mathrm r}(1-p)p^k-\Gamma(1-s)s^k}
     {\kappa_{\mathrm r}p^k-\Gamma s^k}.
\end{equation}
The equal-rate limit is
\begin{equation}\label{eq:equal-rate-survival-hazard-main}
\PP(G\ge k)=p^k(1+\Gamma\Delta\tau\,k),
\qquad
h_k=
\frac{(1-p)+\Gamma\Delta\tau(k(1-p)-p)}
     {1+\Gamma\Delta\tau\,k}.
\end{equation}
In all cases \(h_{k+1}>h_k\).
\end{lemma}

\begin{proof}
Summing \eqref{eq:gap-closed-main} over \(j\ge k\) gives
\eqref{eq:survival-tail-main}.  Since \(g_k=S_k-S_{k+1}\), the hazard has the
equivalent survival-ratio form
\[
h_k=\frac{g_k}{S_k}=1-\frac{S_{k+1}}{S_k}.
\]
Substituting \eqref{eq:survival-tail-main} into this identity gives
\eqref{eq:hazard-main}.  The equal-rate display is the continuous limit.
For the unequal-rate case, subtracting the two survival ratios gives the
standalone identity
\begin{equation}\label{eq:hazard-increment-sign-main}
h_{k+1}-h_k
=
\frac{
  \Gamma\kappa_{\mathrm r}\,
  (p-s)^2\,p^k s^k
}{
  \bigl(\kappa_{\mathrm r}p^k-\Gamma s^k\bigr)
  \bigl(\kappa_{\mathrm r}p^{k+1}-\Gamma s^{k+1}\bigr)
}.
\end{equation}
The numerator in \eqref{eq:hazard-increment-sign-main} is positive because
\(\Gamma,\kappa_{\mathrm r}>0\), \(p,s\in(0,1)\), and the explicit square
\((p-s)^2\) is positive off the equal-rate diagonal.  For the denominator,
\[
\kappa_{\mathrm r}p^m-\Gamma s^m
=(\kappa_{\mathrm r}-\Gamma)S_m,\qquad m=k,k+1,
\]
by \eqref{eq:survival-tail-main}.  Since \(S_m=\PP(G\ge m)>0\), both
denominator factors have sign \(\kappa_{\mathrm r}-\Gamma\), so their product
is positive.  Equivalently, their product is
\((\kappa_{\mathrm r}-\Gamma)^2S_kS_{k+1}>0\).

On the equal-rate diagonal, write
\(y=\Gamma\Delta\tau=\kappa_{\mathrm r}\Delta\tau\) and
\(p=e^{-y}\).  The displayed equal-rate hazard can be rewritten as
\[
h_k
=1-p-\frac{y p}{1+y k}.
\]
Thus
\[
h_{k+1}-h_k
=yp\left(\frac1{1+yk}-\frac1{1+y(k+1)}\right)
=\frac{y^2p}{(1+yk)(1+y(k+1))}>0.
\]
\end{proof}

\begin{proof}[Proof of \Cref{prop:main-record-wrapper-detailed}]
The kernels are those of \Cref{prop:instrument-kernels-main}.  Both nonzero
entries of \(T_1\) land in the \(D\)-column.  Thus, after a click, the next
hidden state is \(D\), independently of the pre-click hidden state and of the
earlier visible record.  Click epochs are therefore regeneration epochs for
the future visible record, and a Palm gap with \(k\) intervening zeros has
probability
\[
g_k=\PP(G=k)=e_D^\top T_0^kT_1\mathbf1,\qquad k\ge0 .
\]
Here
\[
T_0=\begin{pmatrix}p&0\\ b&s\end{pmatrix},
\qquad
T_1\mathbf1=(1-p,a)^\top .
\]
For \(k\ge1\), lower-triangular multiplication gives
\[
T_0^k=
\begin{pmatrix}
p^k&0\\
b\sum_{j=0}^{k-1}s^jp^{k-1-j}&s^k
\end{pmatrix};
\]
for \(k=0\), the same formula is read with the sum empty and \(T_0^0=I\).
Consequently
\[
g_k=b(1-p)\sum_{j=0}^{k-1}s^jp^{k-1-j}+as^k ,
\]
including \(g_0=a\).  If \(p\ne s\), then
\[
\sum_{j=0}^{k-1}s^jp^{k-1-j}=\frac{p^k-s^k}{p-s}.
\]
Using the scalar identities
\[
b=\frac{\kappa_{\mathrm r}(p-s)}{\kappa_{\mathrm r}-\Gamma},
\qquad
a=\frac{\kappa_{\mathrm r}(1-p)-\Gamma(1-s)}
        {\kappa_{\mathrm r}-\Gamma},
\]
one obtains
\[
\begin{aligned}
g_k
&=\frac{\kappa_{\mathrm r}(1-p)(p^k-s^k)}
        {\kappa_{\mathrm r}-\Gamma}
  +\frac{\{\kappa_{\mathrm r}(1-p)-\Gamma(1-s)\}s^k}
        {\kappa_{\mathrm r}-\Gamma} \\
&=\frac{\kappa_{\mathrm r}(1-p)p^k-\Gamma(1-s)s^k}
        {\kappa_{\mathrm r}-\Gamma}.
\end{aligned}
\]
This proves the displayed off-diagonal gap law.  The case \(p=s\), equivalently
\(\Gamma=\kappa_{\mathrm r}\), is not obtained by dividing by \(p-s\); it is
the continuous equal-rate limit recorded in
\eqref{eq:equal-rate-dictionary-main}.  Since the path that remains in \(D\)
for \(k\) zero-labelled steps and then clicks has probability \(s^ka>0\), every
\(g_k\) is positive.  Thus the renewal gap law has full support.

For the hazard, summing the just-proved gap law gives, off the equal-rate
diagonal,
\[
S_k=\PP(G\ge k)=
\frac{\kappa_{\mathrm r}p^k-\Gamma s^k}{\kappa_{\mathrm r}-\Gamma}.
\]
Since \(g_k=S_k-S_{k+1}\),
\[
